\chapter{Producto interno y proyección ortogonal}\label{cap:nuevo14}\label{chap:29}\label{chap:30}
\index{producto interno}\index{Hilbert, espacio de}
\index{proyeccion ortogonal@proyección ortogonal}\index{mejor aproximacion@mejor aproximación}

Una norma mide longitudes; un producto interno permite además medir ángulos y separar componentes perpendiculares. Esta información adicional convierte ciertos problemas de aproximación en problemas geométricos. En geometría euclidiana, el punto de una recta más cercano a un punto exterior se obtiene trazando una perpendicular. El mismo principio funciona en dimensión infinita, pero la existencia del punto más cercano exige atención.

Construiremos ese puente desde los axiomas del producto interno y la desigualdad de Cauchy--Schwarz hasta la proyección ortogonal. Incluiremos los espacios de sucesiones con un conjunto de índices no necesariamente numerable y distinguiremos la \emph{caracterización} del mejor aproximante de su \emph{existencia}, para identificar exactamente dónde interviene la completitud. Se presuponen las propiedades de las normas, la convergencia de series numéricas y la completitud. El capítulo reúne los temas de las lecturas 29 y 30 de \cite{valencia}, con las correcciones de hipótesis desarrolladas en sus ejercicios.

\section{Producto interno, norma inducida y completitud}
\label{lectura:29}
\begin{definicion}
Sea $V$ un espacio vectorial sobre $\F=\R$ o $\C$. Un \emph{producto interno} es una función $\ip{\cdot}{\cdot}:V\times V\to\F$ que cumple
\[
 \ip{x}{x}\geq0,\qquad \ip{x}{x}=0\Longleftrightarrow x=0,
\]
\[
 \ip{\alpha x+\beta y}{z}=\alpha\ip{x}{z}+\beta\ip{y}{z},
 \qquad \ip{x}{y}=\overline{\ip{y}{x}}.
\]
La última propiedad se llama \emph{simetría conjugada}. En todo este volumen el producto interno es lineal en la primera entrada.
\end{definicion}
La segunda entrada es antilineal, pues
\[
 \ip{x}{\alpha y+\beta z}
 =\overline{\ip{\alpha y+\beta z}{x}}
 =\bar\alpha\ip{x}{y}+\bar\beta\ip{x}{z}.
\]
Sobre $\R$ no aparecen conjugados. Es útil escribirlos siempre al principio: muchos errores sobre adjuntos proceden de cambiar inadvertidamente la convención.

\begin{teorema}[Cauchy--Schwarz]\label{ch29:cs}
\index{Cauchy--Schwarz, desigualdad de}
Si $\norm{x}=\sqrt{\ip{x}{x}}$, entonces
\[
 |\ip{x}{y}|\leq\norm{x}\norm{y}.
\]
\end{teorema}
\begin{proof}
Si $y=0$, ambos miembros son cero. Si $y\ne0$, para cualquier $\alpha\in\F$,
\[
0\leq\norm{x-\alpha y}^{2}
=\norm{x}^{2}-\bar\alpha\ip{x}{y}-\alpha\ip{y}{x}
 +|\alpha|^{2}\norm{y}^{2}.
\]
La elección $\alpha=\ip{x}{y}/\norm{y}^{2}$ cancela la componente de $x$ paralela a $y$ y da
\[
0\leq\norm{x}^{2}-\frac{|\ip{x}{y}|^{2}}{\norm{y}^{2}}.
\]
Multiplicar por $\norm{y}^{2}$ y tomar raíces produce la desigualdad.
\end{proof}

\begin{proposicion}\label{ch29:norma}
\index{paralelogramo, identidad del}
La función $x\mapsto\sqrt{\ip{x}{x}}$ es una norma. Además,
\begin{align}
 \norm{x+y}^{2}&=\norm{x}^{2}+2\operatorname{Re}\ip{x}{y}+\norm{y}^{2},\label{ch29:expansion}\\
 \norm{x+y}^{2}+\norm{x-y}^{2}&=2\norm{x}^{2}+2\norm{y}^{2}.\label{ch29:paralelogramo}
\end{align}
\end{proposicion}
\begin{proof}
La positividad y la separación del cero son axiomas. La sesquilinealidad da $\ip{\alpha x}{\alpha x}=|\alpha|^{2}\ip{x}{x}$ y, por tanto, homogeneidad absoluta de la norma. Desarrollar $\ip{x+y}{x+y}$ demuestra \eqref{ch29:expansion}; Cauchy--Schwarz implica
\[
 \norm{x+y}^{2}\leq \norm{x}^{2}+2\norm{x}\norm{y}+\norm{y}^{2}
 =(\norm{x}+\norm{y})^{2}.
\]
La raíz cuadrada preserva el orden de números no negativos, de modo que se obtiene la desigualdad triangular. Sumar las expansiones correspondientes a $x+y$ y $x-y$ cancela los términos cruzados y prueba \eqref{ch29:paralelogramo}.
\end{proof}

\begin{definicion}
Un espacio con producto interno es un \emph{espacio de Hilbert} cuando es completo para la norma inducida. Sin completitud se habla también de espacio \emph{prehilbertiano}.
\end{definicion}

\begin{ejemplo}[Tres modelos]
En $\F^{n}$ se utiliza $\ip{x}{y}=\sum_{j=1}^{n}x_j\bar y_j$. Es un espacio de Hilbert porque una sucesión de Cauchy converge coordenada a coordenada y, al haber solamente $n$ coordenadas, converge en norma.

En $\ell^{2}$, el producto $\ip{x}{y}=\sum_{j\geq1}x_j\bar y_j$ está bien definido: Cauchy--Schwarz para sumas finitas acota $\sum_{j=1}^{N}|x_jy_j|$ por $\norm{x}_2\norm{y}_2$. La completitud se demostrará, en mayor generalidad, en el ejercicio~\ref{ej:lectura-29-8}.

Si $a<b$, en $C([a,b])$ se tiene
\[
 \ip{f}{g}=\frac1{b-a}\int_a^b f(t)\overline{g(t)}\,dt.
\]
Si $f$ no es cero, existe un intervalo de longitud positiva donde $|f|$ está acotado inferiormente por una constante positiva. Por ello $\int|f|^{2}>0$. Sin embargo, este espacio no es completo. En $[0,1]$ considérese $f_n$ igual a cero hasta $1/2$, lineal de cero a uno entre $1/2$ y $1/2+1/n$, e igual a uno después. Si $s=\mathbf1_{(1/2,1]}$, entonces $\int|f_n-s|^{2}\leq1/n$, por lo que $(f_n)$ es de Cauchy. Si convergiera en la norma integral a una función continua $f$, la desigualdad triangular daría $\int|f-s|^{2}=0$. La continuidad obligaría a $f=0$ en $(0,1/2)$ y a $f=1$ en $(1/2,1)$, incompatible con la continuidad en $1/2$. Un cambio afín de variable traslada el ejemplo a cualquier $[a,b]$.
\end{ejemplo}

\section{Ortogonalidad, sumas y desigualdad de Bessel}
\index{ortogonalidad}\index{Bessel, desigualdad de}
Se escribe $x\perp y$ si $\ip{x}{y}=0$. Para $S\subset V$ se define
\[
 S^{\perp}=\{x\in V:\ip{x}{s}=0\text{ para todo }s\in S\}.
\]
La linealidad demuestra que $S^{\perp}$ es un subespacio. Es cerrado: si $x_n\in S^{\perp}$ y $x_n\to x$, entonces, para cada $s\in S$,
$|\ip{x}{s}|=|\ip{x-x_n}{s}|\leq\norm{x-x_n}\norm{s}\to0$.
El único vector perpendicular a todos los vectores de $V$ es cero, porque es entonces perpendicular a sí mismo. Dos vectores ortogonales \emph{no nulos} son linealmente independientes: de $\alpha x+\beta y=0$ se obtiene $\alpha\norm{x}^{2}=0$ y $\beta\norm{y}^{2}=0$. La condición de no nulidad es indispensable.

\begin{teorema}[Pitágoras]\label{ch29:pitagoras}
\index{Pitagoras@Pitágoras, teorema de}
Si $x_1,\dots,x_n$ son ortogonales a pares, entonces
\[
 \Big\|\sum_{j=1}^{n}x_j\Big\|^{2}=\sum_{j=1}^{n}\norm{x_j}^{2}.
\]
\end{teorema}
\begin{proof}
Por sesquilinealidad, el miembro izquierdo es $\sum_{j,k}\ip{x_j}{x_k}$. Los términos con $j\ne k$ son cero; los restantes son $\ip{x_j}{x_j}=\norm{x_j}^{2}$. Esta expansión también explica por qué la suma de varios vectores perpendiculares a $y$ sigue siendo perpendicular a $y$.
\end{proof}

\begin{definicion}\label{ch29:suma}
Para una familia $(a_s)_{s\in S}$ de números no negativos se define
\[
 \sum_{s\in S}a_s=\sup\Big\{\sum_{s\in F}a_s:F\subset S\text{ finito}\Big\}.
\]
Cada índice se cuenta una sola vez. Un conjunto $E\subset V$ es \emph{ortonormal} si sus vectores tienen norma uno y son ortogonales dos a dos.
\end{definicion}
\index{suma!familia no numerable}\index{conjunto ortonormal}

\begin{lema}\label{ch29:soporte}
Si $A=\sum_{s\in S}a_s<\infty$, entonces $\{s:a_s>0\}$ es a lo sumo numerable. Cualquier enumeración sin repeticiones de ese conjunto tiene suma $A$.
\end{lema}
\begin{proof}
Para cada $m\geq1$, el conjunto $S_m=\{s:a_s>1/m\}$ es finito: si contuviera $r$ elementos distintos, $r/m<\sum_{s\in F}a_s\leq A$; no pueden existir subconjuntos finitos de cardinal arbitrariamente grande. Como todo número positivo supera $1/m$ para algún $m$, el soporte es $\bigcup_m S_m$. Si lo enumeramos como $s_1,s_2,\dots$, las sumas parciales crecen y están acotadas por $A$, de modo que tienen límite $L\leq A$. Para cualquier conjunto finito $F$, sus índices con coeficiente no nulo aparecen entre los primeros $N$ términos para algún $N$, y entonces $\sum_{s\in F}a_s\leq L$. Tomar el supremo da $A\leq L$. El soporte finito o vacío se trata mediante una suma finita o la suma cero.
\end{proof}

\begin{teorema}[Bessel]\label{ch29:bessel}
Si $E$ es ortonormal, entonces para todo $x\in V$,
\[
 \sum_{e\in E}|\ip{x}{e}|^{2}\leq\norm{x}^{2}.
\]
\end{teorema}
\begin{proof}
Para $F\subset E$ finito, sea $p_F=\sum_{e\in F}\ip{x}{e}e$. Para cada $u\in F$,
\[
 \ip{x-p_F}{u}=\ip{x}{u}-\sum_{e\in F}\ip{x}{e}\ip{e}{u}=0.
\]
Por tanto, $x-p_F\perp p_F$ y Pitágoras da
\[
 \norm{x}^{2}=\norm{x-p_F}^{2}+\norm{p_F}^{2}
 =\norm{x-p_F}^{2}+\sum_{e\in F}|\ip{x}{e}|^{2}.
\]
Tomamos el supremo sobre $F$. El lema anterior demuestra que, para este $x$ fijo, solamente una cantidad a lo sumo numerable de coeficientes es distinta de cero. No demuestra que $E$ mismo sea numerable; en $\ell^{2}(S)$ puede haber un vector coordenado por cada elemento de un conjunto no numerable $S$.
\end{proof}

\section{Ejercicios de producto interno, polarización y completitud}

\begin{ejercicio}[Lectura 29, ejercicio 1]\label{ej:lectura-29-1}
Caracterizar la igualdad en Cauchy--Schwarz.
\end{ejercicio}


\begin{ejercicio}[Lectura 29, ejercicio 2: corrección del recíproco]\label{ej:lectura-29-2}
Analizar si $\norm{v_1+\cdots+v_n}^{2}=\sum_j\norm{v_j}^{2}$ implica ortogonalidad por pares en un espacio real.
\end{ejercicio}


\begin{ejercicio}[Lectura 29, ejercicio 3: identidad del paralelogramo]\label{ej:lectura-29-3}
Demostrar que una norma procede de un producto interno si y solo si
\[
 \norm{x+y}^{2}+\norm{x-y}^{2}=2\norm{x}^{2}+2\norm{y}^{2}.
\]
Aplicarlo a (i) $\ell^{1}$ y (ii) la norma uniforme en $B(S,W)$.
\end{ejercicio}


\begin{ejercicio}[Lectura 29, ejercicio 4]\label{ej:lectura-29-4}
Demostrar la continuidad conjunta del producto interno.
\end{ejercicio}


\begin{ejercicio}[Lectura 29, ejercicio 5: identidad polar]\label{ej:lectura-29-5}
\index{polarizacion@polarización, identidad de}
Si $\zeta$ es una raíz primitiva de la unidad de orden $n\geq3$, probar
\[
 \ip{x}{y}=\frac1n\sum_{k=0}^{n-1}\zeta^k\norm{x+\zeta^ky}^{2}.
\]
\end{ejercicio}


\begin{ejercicio}[Lectura 29, ejercicio 6]\label{ej:lectura-29-6}
Sea $h:[a,b]\to\F$ continua y sin ceros, con $a<b$. Determinar cuándo $\ip{f}{g}_h=\int_a^b f\bar g h$ es producto interno, y discutir el caso continuo a trozos.
\end{ejercicio}


\begin{ejercicio}[Lectura 29, ejercicio 7]\label{ej:lectura-29-7}
\index{ell2@\ensuremath{\ell^2}!ponderado}
Para $w_n>0$, sea $\ell^{2}_w=\{x:\sum_n w_n|x_n|^{2}<\infty\}$. Probar (i) que es espacio vectorial, (ii) que $\ip{x}{y}_w=\sum_nw_nx_n\bar y_n$ es producto interno y (iii) que es completo.
\end{ejercicio}


\begin{ejercicio}[Lectura 29, ejercicio 8]\label{ej:lectura-29-8}
\index{ell2@\ensuremath{\ell^2}!indices arbitrarios@índices arbitrarios}
Para cualquier conjunto $S$, definir $\ell^{2}(S)$ por $\sum_{s\in S}|f(s)|^{2}<\infty$, con la suma de la definición \ref{ch29:suma}. Probar (i) la estructura vectorial, (ii) la existencia del producto $\ip{f}{g}=\sum_{s\in S}f(s)\overline{g(s)}$ y (iii) la completitud.
\end{ejercicio}


\section{Caracterización geométrica de la mejor aproximación}
\label{lectura:30}
La identidad del paralelogramo caracteriza las normas que admiten ángulos, y Bessel cuantifica cuánta longitud se captura al conservar componentes ortogonales. Convertiremos ahora esta idea en un procedimiento de mejor aproximación: la proyección ortogonal. Primero probaremos qué propiedad debe satisfacer un mejor aproximante, cuando existe; después mostraremos que la completitud permite construirlo.
\begin{teorema}\label{ch30:caracterizacion}
Sea $V$ un espacio con producto interno y $M\subset V$ un subespacio. Para $x\in V$ y $m\in M$ son equivalentes:
\begin{enumerate}[label=(\roman*)]
 \item $x-m\in M^{\perp}$;
 \item $\norm{x-m}=\dist(x,M)$.
\end{enumerate}
Si se cumplen, $m$ es único y para todo $u\in M$,
\[
 \norm{x-u}^{2}=\norm{x-m}^{2}+\norm{m-u}^{2}.
\]
En particular, $\dist(x,M)^{2}=\norm{x}^{2}-\norm{m}^{2}$.
\end{teorema}
\begin{proof}
Si $x-m\perp M$, entonces $x-m\perp m-u$, de modo que Pitágoras da la igualdad anunciada. Para $u\ne m$, el segundo sumando es positivo; por ello $m$ es el único minimizador. Con $u=0$ obtenemos la fórmula de la distancia.

Recíprocamente, supongamos que $m$ minimiza. Fijemos $u\in M\setminus\{0\}$ y sea $r=x-m$. Como $m+\alpha u\in M$ para cada $\alpha\in\F$,
\[
0\leq\norm{r-\alpha u}^{2}-\norm{r}^{2}
=-2\operatorname{Re}(\bar\alpha\ip{r}{u})+|\alpha|^{2}\norm{u}^{2}.
\]
Elegimos $\alpha=\ip{r}{u}/\norm{u}^{2}$. El lado derecho se vuelve $-|\ip{r}{u}|^{2}/\norm{u}^{2}$. Debe ser no negativo, por lo que $\ip{r}{u}=0$. Para $u=0$ la igualdad ya se cumple. Así $r\in M^{\perp}$.
\end{proof}

\begin{definicion}
El vector del teorema se llama \emph{proyección ortogonal} de $x$ sobre $M$ y se escribe $P_Mx$. Un subespacio es \emph{ortogonalmente complementado}, o \emph{distinguido} en las lecturas originales, cuando dicha proyección existe para todo $x\in V$.
\end{definicion}
\index{subespacio!distinguido}\index{subespacio!ortogonalmente complementado}

\begin{ejemplo}[La distancia puede no alcanzarse]
En $\ell^{2}$, sea $M=c_{00}=\spanop\{e_n:n\geq1\}$. Los truncamientos de cualquier sucesión convergen en norma, de modo que $M$ es denso. Para $x=(1/n)_{n\geq1}$ tenemos $\dist(x,M)=0$, pero $x\notin M$. Si $u\in M$ y $u_n=0$ para $n\geq N$, entonces
\[
 \norm{x-u}^{2}\geq\sum_{n=N}^{\infty}\frac1{n^{2}}>0.
\]
Por ello ningún $u$ alcanza la distancia. El obstáculo es que el límite de una sucesión de aproximantes puede salir del subespacio.
\end{ejemplo}

\section{Existencia y descomposición ortogonal}
\begin{teorema}[Teorema de proyección]\label{ch30:proyeccion}
En un espacio con producto interno:
\begin{enumerate}[label=(\roman*)]
 \item todo subespacio ortogonalmente complementado es cerrado;
 \item todo subespacio completo es ortogonalmente complementado.
\end{enumerate}
Por tanto, en un espacio de Hilbert, un subespacio admite proyección ortogonal si y solo si es cerrado.
\end{teorema}
\begin{proof}
\emph{(i)} Si $x\in\cl M$, entonces $\dist(x,M)=0$. Si $m=P_Mx$ existe, $\norm{x-m}=0$, luego $x=m\in M$.

\emph{(ii)} Fijemos $x$ y escribamos $d=\dist(x,M)$. Por definición de ínfimo, para cada $n$ se puede elegir $m_n\in M$ con $d\leq\norm{x-m_n}<d+1/n$. El paralelogramo aplicado a $x-m_n$ y $x-m_k$ da
\[
 \norm{m_n-m_k}^{2}
 =2\norm{x-m_n}^{2}+2\norm{x-m_k}^{2}
 -4\left\|x-\frac{m_n+m_k}{2}\right\|^{2}.
\]
El punto medio pertenece a $M$, de modo que su distancia a $x$ es al menos $d$. Por ello
\[
0\leq\norm{m_n-m_k}^{2}
\leq2\norm{x-m_n}^{2}+2\norm{x-m_k}^{2}-4d^{2}\longrightarrow0.
\]
Esto prueba que $(m_n)$ es de Cauchy. Por completitud de $M$, converge a $m\in M$, y la continuidad de la norma implica $\norm{x-m}=d$. El teorema \ref{ch30:caracterizacion} termina la prueba.

Si el espacio ambiente es Hilbert, sus subespacios cerrados son completos: el límite ambiente de una sucesión de Cauchy permanece en el subespacio por cerradura. Esto explica la última equivalencia.
\end{proof}

\begin{teorema}\label{ch30:descomposicion}
\index{descomposicion ortogonal@descomposición ortogonal}
Un subespacio $M$ es ortogonalmente complementado si y solo si
\[
 V=M\oplus M^{\perp}.
\]
En ese caso $(M^{\perp})^{\perp}=M$, y $M^{\perp}$ también es ortogonalmente complementado.
\end{teorema}
\begin{proof}
Si existe $P_Mx$, escribimos $x=P_Mx+(x-P_Mx)$, con el primer sumando en $M$ y el segundo en $M^{\perp}$. Además $M\cap M^{\perp}=\{0\}$, pues un vector de la intersección es perpendicular a sí mismo. La descomposición es, por tanto, única.

Recíprocamente, si $x=m+n$ con $m\in M$ y $n\in M^{\perp}$, el teorema \ref{ch30:caracterizacion} demuestra que $m=P_Mx$. Es el componente en $M$, no el componente perpendicular.

Siempre $M\subset(M^{\perp})^{\perp}$. Si $x\in(M^{\perp})^{\perp}$, descompongamos $x=m+n$. Entonces $0=\ip{x}{n}=\ip{m}{n}+\ip{n}{n}=\norm{n}^{2}$, y $x=m\in M$. Así $V=M^{\perp}\oplus(M^{\perp})^{\perp}$; la equivalencia ya demostrada prueba que $M^{\perp}$ es complementado.
\end{proof}

\begin{corolario}\label{ch30:doble}
\index{complemento ortogonal!doble}
Si $H$ es Hilbert y $S\subset H$, entonces
\[
 (S^{\perp})^{\perp}=\cl{\spanop S}.
\]
\end{corolario}
\begin{proof}
Un vector perpendicular a $S$ es perpendicular a todas sus combinaciones lineales y a sus límites, por continuidad del producto interno. Por tanto, con $M=\cl{\spanop S}$ se tiene $S^{\perp}=M^{\perp}$. Como $M$ es cerrado en Hilbert, es complementado y el teorema anterior da $(M^{\perp})^{\perp}=M$.
\end{proof}

\section{Calcular una proyección sin ortonormalizar}\label{ampl:gram}
El teorema de proyección garantiza un mejor aproximante, pero no siempre lo expresa en las coordenadas que tenemos disponibles. Si un subespacio está generado por monomios, datos experimentales u otras funciones sencillas, conviene poder trabajar directamente con esa lista. La matriz de Gram convierte la condición de ortogonalidad del error en un sistema de ecuaciones lineales.

Esta ampliación utiliza \cite{max-gram,max-proyeccion} como orientación. Las demostraciones se escriben con la convención del volumen: el producto interno es lineal en el primer argumento. Las fórmulas complejas dependen de esa convención; no basta con copiar una matriz sin comprobar dónde aparecen los conjugados.
\index{Gram, matriz de}\index{ecuaciones normales}\index{minimos cuadrados@mínimos cuadrados}

\begin{definicion}
Para una lista $a_1,\ldots,a_m$ de vectores de un espacio con producto interno, su \emph{matriz de Gram} será
\[
 G=(G_{ij})_{i,j=1}^m,\qquad G_{ij}=\ip{a_j}{a_i}.
\]
Para una matriz compleja, $A^*$ denota su transpuesta conjugada. Los vectores de coeficientes se escriben como columnas.
\end{definicion}

\begin{proposicion}\label{ampl:gram-positiva}
La matriz $G$ es hermitiana y semidefinida positiva. Es definida positiva, y por tanto invertible, si y solo si $a_1,\ldots,a_m$ son linealmente independientes. Para cada columna $c=(c_j)$,
\[
 c^*Gc=\left\|\sum_{j=1}^mc_ja_j\right\|^2.
\]
\end{proposicion}
\begin{proof}
La simetría conjugada da $\overline{G_{ji}}=\overline{\ip{a_i}{a_j}}=\ip{a_j}{a_i}=G_{ij}$, luego $G^*=G$. Además,
\[
 c^*Gc=\sum_{i,j}\overline{c_i}c_j\ip{a_j}{a_i}
 =\ip{\sum_jc_ja_j}{\sum_ic_ia_i}\ge0.
\]
Si la lista es independiente y $c\ne0$, el vector dentro de la norma no es cero, por lo que la expresión es positiva. Si la lista es dependiente, existe $c\ne0$ que produce el vector cero, y la expresión no es estrictamente positiva. Finalmente, para una lista independiente, $Gc=0$ implicaría $c^*Gc=0$ y por tanto $c=0$; una matriz cuadrada con núcleo cero es invertible.
\end{proof}

\begin{teorema}[Ecuaciones normales de la proyección]\label{ampl:gram-proyeccion}
Sean $a_1,\ldots,a_m$ linealmente independientes y $M=\spanop\{a_1,\ldots,a_m\}$. Para $u$ en el espacio ambiente defina la columna $\alpha$ por $\alpha_i=\ip{u}{a_i}$. El sistema
\[
 Gc=\alpha
\]
tiene solución única, y
\[
 P_Mu=\sum_{j=1}^mc_ja_j,\qquad
 \dist(u,M)^2=\norm u^2-\alpha^*G^{-1}\alpha.
\]
No se requiere que el espacio ambiente sea completo.
\end{teorema}
\begin{proof}
Por la proposición anterior existe $c=G^{-1}\alpha$. Sea $p=\sum_jc_ja_j$. Para cada $i$,
\[
 \ip{u-p}{a_i}=\alpha_i-\sum_jc_j\ip{a_j}{a_i}
 =\alpha_i-(Gc)_i=0.
\]
Si $v=\sum_id_ia_i\in M$, entonces
$\ip{u-p}{v}=\sum_i\overline{d_i}\ip{u-p}{a_i}=0$. El residuo es ortogonal a todo $M$, de modo que el teorema~\ref{ch30:caracterizacion} prueba que $p$ es el único mejor aproximante. El argumento sólo usó sumas finitas y un sistema invertible, por lo que no necesita completitud ambiente.

Como $G$ es hermitiana, también lo es $G^{-1}$: al tomar adjuntos en $GG^{-1}=I$ se obtiene $(G^{-1})^*G=I$. Por ello,
\[
 \norm p^2=c^*Gc=(G^{-1}\alpha)^*G(G^{-1}\alpha)
 =\alpha^*G^{-1}\alpha.
\]
La descomposición $u=p+(u-p)$ es ortogonal. Pitágoras da la fórmula del error.
\end{proof}

En la práctica se resuelve $Gc=\alpha$; escribir $G^{-1}$ en una fórmula no obliga a calcular todas las entradas de la inversa. Cuando la lista ya es ortonormal, $G=I$ y recuperamos $c_i=\ip{u}{a_i}$. La matriz de Gram registra precisamente lo que cambia cuando las direcciones no son ortogonales.

\begin{ejemplo}[Matrices y ajuste por mínimos cuadrados]
En $\F^n$ con el producto interno usual, reúna los vectores $a_j$ como columnas de una matriz $A$ de tamaño $n\times m$. Entonces
\[
 G=A^*A,\qquad \alpha=A^*u.
\]
En efecto, la entrada $(i,j)$ de $A^*A$ es $\sum_k\overline{(a_i)_k}(a_j)_k=\ip{a_j}{a_i}$. Las ecuaciones normales son $A^*Ac=A^*u$ y minimizan $\norm{u-Ac}_2^2$. Si las columnas son independientes, la proyección queda expresada por
\[
 P_Mu=A(A^*A)^{-1}A^*u.
\]
La fórmula es una aplicación finita del teorema anterior; no presupone todavía la teoría del adjunto de un operador entre espacios de Hilbert.
\end{ejemplo}

\section{Ejercicios sobre Gram y mejores aproximaciones}
Los ejercicios siguientes son elaboraciones propias de esta edición. El último desarrolla también el problema sobre rango de Gram propuesto en \cite{max-gram}; las soluciones y las aplicaciones concretas se presentan aquí completas.

\begin{ejercicio}[Ampliación: aproximar una parábola por una recta]\label{ej:ampl-gram-1}
En $C([0,1],\R)$ con $\ip fg=\int_0^1f(t)g(t)\,dt$, encuentre el polinomio $a+bt$ que minimiza $\int_0^1(t^2-a-bt)^2\,dt$. Calcule el error mínimo y verifique directamente la ortogonalidad del residuo.
\end{ejercicio}


\begin{ejercicio}[Ampliación: dónde va el conjugado]\label{ej:ampl-gram-2}
En $\C^2$, con $\ip xy=x_1\overline{y_1}+x_2\overline{y_2}$, proyecte $u=(1,1)$ sobre la recta compleja generada por $a=(1,i)$. Compruebe la norma del error y explique por qué usar $\ip au/\norm a^2$ como coeficiente daría una respuesta incorrecta.
\end{ejercicio}


\begin{ejercicio}[Ampliación: tres datos y dos parámetros]\label{ej:ampl-gram-3}
Ajuste una recta $y=a+bt$ a los datos $(0,1)$, $(1,2)$ y $(2,2)$ minimizando la suma de los cuadrados de los tres errores. Calcule los parámetros, los valores ajustados y el error mínimo. Explique por qué el problema tiene una solución única aunque ninguna recta pase por los tres datos.
\end{ejercicio}


\begin{ejercicio}[Ampliación: cambiar de base sin cambiar la proyección]\label{ej:ampl-gram-4}
Sea $a_1,\ldots,a_m$ una lista independiente y $C=(C_{ij})$ una matriz invertible. Defina $b_j=\sum_i a_iC_{ij}$. Pruebe que las matrices de Gram y los datos de proyección satisfacen $G_b=C^*G_aC$ y $\alpha_b=C^*\alpha_a$, y relacione las coordenadas del mismo vector proyectado. Ilustre el resultado con $a_1=(1,0)$, $a_2=(1,\eps)$, $\eps>0$, y $u=(0,1)$ en $\R^2$.
\end{ejercicio}


\begin{ejercicio}[Ampliación: listas dependientes y coeficientes no únicos]\label{ej:ampl-gram-5}
Permita ahora que $a_1,\ldots,a_m$ sean dependientes. Defina $A:\F^m\to V$ por $Ac=\sum_jc_ja_j$. Pruebe que $\ker G=\ker A$, que $\operatorname{rango}G=\dim\spanop\{a_j\}$ y que $Gc=\alpha$ sigue teniendo soluciones para $\alpha_i=\ip u{a_i}$. Muestre que todas producen el mismo mejor aproximante. Calcule el sistema para $a_1=1$, $a_2=t$, $a_3=1+t$ y $u=t^2$, con producto integral en $[0,1]$.
\end{ejercicio}


La fórmula de Gram concreta la proyección en subespacios finitos. Los ejercicios que siguen vuelven a la geometría general, donde intervienen conjuntos convexos, complementos ortogonales y familias que pueden ser infinitas.

\section{Ejercicios de proyección y complementos ortogonales}

\begin{ejercicio}[Lectura 30, ejercicio 1]\label{ej:lectura-30-1}
Sea $C$ un subconjunto convexo, completo y no vacío de un espacio con producto interno. Probar que para cada $x$ existe un único $c\in C$ tal que $\norm{x-c}=\dist(x,C)$.
\end{ejercicio}


\begin{ejercicio}[Lectura 30, ejercicio 2: Parseval sobre un subespacio]\label{ej:lectura-30-2}
Sea $M$ ortogonalmente complementado y $E\subset M$ ortonormal con $\spanop E$ denso en $M$. Para $x\in V$ y $m=P_Mx$, probar
\[
 \norm{m}^{2}=\sum_{e\in E}|\ip{x}{e}|^{2}.
\]
\end{ejercicio}


\begin{ejercicio}[Lectura 30, ejercicio 3: tres perpendiculares]\label{ej:lectura-30-3}
Sean $M\subset N$ subespacios ortogonalmente complementados. Probar $P_MP_Nx=P_Mx$.
\end{ejercicio}


\begin{ejercicio}[Lectura 30, ejercicio 4: complementos ortogonales]\label{ej:lectura-30-4}
Para subespacios $M,N$, estudiar
\begin{enumerate}[label=(\roman*)]
 \item $(M+N)^{\perp}=M^{\perp}\cap N^{\perp}$;
 \item $(M\cap N)^{\perp}=\cl{M^{\perp}+N^{\perp}}$;
 \item $M\subset(M^{\perp})^{\perp}$.
\end{enumerate}
En (ii) suponemos que $M,N$ son subespacios cerrados de un espacio de Hilbert. Explicaremos por qué no basta, en general, suponer que $M,N$ sean distinguidos en un espacio incompleto. Los incisos (i) y (iii) no requieren cerradura ni completitud.
\end{ejercicio}


\begin{ejercicio}[Lectura 30, ejercicio 5]\label{ej:lectura-30-5}
Probar que $P_M$ es lineal, continuo, idempotente y simétrico respecto al producto interno. Determinar su norma.
\end{ejercicio}


\begin{ejercicio}[Lectura 30, ejercicio 6]\label{ej:lectura-30-6}
Sea $P:V\to V$ lineal, $P^{2}=P$, y $\ip{Px}{y}=\ip{x}{Py}$ para todos $x,y$. Probar que $P$ es una proyección ortogonal.
\end{ejercicio}


\section{La descomposición como control del error}
La fórmula $x=P_Mx+(x-P_Mx)$ separa la información representable en $M$ del error perpendicular. La perpendicularidad no es solamente una imagen geométrica: demuestra unicidad, controla la norma del error y convierte la aproximación por un sistema ortonormal en una suma de cuadrados. En el capítulo~\ref{cap:nuevo15} haremos esa suma con todos los vectores de una base ortonormal.
